\documentclass[10pt,a4paper]{amsart}
\usepackage{amsmath,amssymb,amsthm,graphicx}
\usepackage{verbatim}
\newenvironment{reference}{\comment}{\endcomment}
%\newenvironment{reference}{}{}
\newenvironment{slogan}{\comment}{\endcomment}
\newenvironment{history}{\comment}{\endcomment}

% For commutative diagrams we use Xy-pic
\usepackage[all]{xy}

% We use 2cell for 2-commutative diagrams.
\xyoption{2cell}
% Theorem environments.
%
\theoremstyle{plain}
\newtheorem{theorem}[subsection]{Theorem}
\newtheorem{proposition}[subsection]{Proposition}
\newtheorem{lemma}[subsection]{Lemma}

\theoremstyle{definition}
\newtheorem{definition}[subsection]{Definition}
\newtheorem{example}[subsection]{Example}
\newtheorem{exercise}[subsection]{Exercise}
\newtheorem{situation}[subsection]{Situation}

\theoremstyle{remark}
\newtheorem{remark}[subsection]{Remark}
\newtheorem{remarks}[subsection]{Remarks}

\numberwithin{equation}{subsection}

% Macros
%
\def\lim{\mathop{\mathrm{lim}}\nolimits}
\def\colim{\mathop{\mathrm{colim}}\nolimits}
\def\Spec{\mathop{\mathrm{Spec}}}
\def\Hom{\mathop{\mathrm{Hom}}\nolimits}
\def\Ext{\mathop{\mathrm{Ext}}\nolimits}
\def\SheafHom{\mathop{\mathcal{H}\!\mathit{om}}\nolimits}
\def\SheafExt{\mathop{\mathcal{E}\!\mathit{xt}}\nolimits}
\def\Sch{\mathit{Sch}}
\def\Mor{\mathop{\mathrm{Mor}}\nolimits}
\def\Ob{\mathop{\mathrm{Ob}}\nolimits}
\def\Sh{\mathop{\mathit{Sh}}\nolimits}
\def\NL{\mathop{N\!L}\nolimits}
\def\CH{\mathop{\mathrm{CH}}\nolimits}
\def\proetale{{pro\text{-}\acute{e}tale}}
\def\etale{{\acute{e}tale}}
\def\QCoh{\mathit{QCoh}}
\def\Ker{\mathop{\mathrm{Ker}}}
\def\Im{\mathop{\mathrm{Im}}}
\def\Coker{\mathop{\mathrm{Coker}}}
\def\Coim{\mathop{\mathrm{Coim}}}

% Boxtimes
%
\DeclareMathSymbol{\boxtimes}{\mathbin}{AMSa}{"02}

%
% Macros for moduli stacks/spaces
%
\def\QCohstack{\mathcal{QC}\!\mathit{oh}}
\def\Cohstack{\mathcal{C}\!\mathit{oh}}
\def\Spacesstack{\mathcal{S}\!\mathit{paces}}
\def\Quotfunctor{\mathrm{Quot}}
\def\Hilbfunctor{\mathrm{Hilb}}
\def\Curvesstack{\mathcal{C}\!\mathit{urves}}
\def\Polarizedstack{\mathcal{P}\!\mathit{olarized}}
\def\Complexesstack{\mathcal{C}\!\mathit{omplexes}}
% \Pic is the operator that assigns to X its picard group, usage \Pic(X)
% \Picardstack_{X/B} denotes the Picard stack of X over B
% \Picardfunctor_{X/B} denotes the Picard functor of X over B
\def\Pic{\mathop{\mathrm{Pic}}\nolimits}
\def\Picardstack{\mathcal{P}\!\mathit{ic}}
\def\Picardfunctor{\mathrm{Pic}}
\def\Deformationcategory{\mathcal{D}\!\mathit{ef}}
\usepackage{cleveref}
\crefname{originalreference}{original reference}{original references}
\Crefname{originalreference}{Original reference}{Original references}
\makeatletter
\newlabel{coherent-lemma-coherent-Noetherian}{{01XZ}{source}{}{sourcepointer-1}{}}
\newlabel{coherent-lemma-coherent-Noetherian@cref}{{[originalreference][0][]01XZ}{[1][0][]source}{}{}{}}
\makeatother
\crefname{originalreference}{original reference}{original references}
\Crefname{originalreference}{Original reference}{Original references}
\makeatletter
\newlabel{perfect-lemma-quasi-coherence-tensor-product}{{08DX}{source}{}{sourcepointer-2}{}}
\newlabel{perfect-lemma-quasi-coherence-tensor-product@cref}{{[originalreference][0][]08DX}{[1][0][]source}{}{}{}}
\makeatother
\crefname{originalreference}{original reference}{original references}
\Crefname{originalreference}{Original reference}{Original references}
\makeatletter
\newlabel{algebra-lemma-tor-noetherian}{{0AZ4}{source}{}{sourcepointer-3}{}}
\newlabel{algebra-lemma-tor-noetherian@cref}{{[originalreference][0][]0AZ4}{[1][0][]source}{}{}{}}
\makeatother
\crefname{originalreference}{original reference}{original references}
\Crefname{originalreference}{Original reference}{Original references}
\makeatletter
\newlabel{perfect-lemma-identify-pseudo-coherent-noetherian}{{08E8}{source}{}{sourcepointer-4}{}}
\newlabel{perfect-lemma-identify-pseudo-coherent-noetherian@cref}{{[originalreference][0][]08E8}{[1][0][]source}{}{}{}}
\makeatother
\crefname{originalreference}{original reference}{original references}
\Crefname{originalreference}{Original reference}{Original references}
\makeatletter
\newlabel{cohomology-lemma-tensor-pseudo-coherent}{{09J3}{source}{}{sourcepointer-5}{}}
\newlabel{cohomology-lemma-tensor-pseudo-coherent@cref}{{[originalreference][0][]09J3}{[1][0][]source}{}{}{}}
\makeatother
\title{Coherence of Tor sheaves and derived tensor products on locally Noetherian schemes}
\author{The Stacks Project Authors}
\date{stacks-project:intersection.tex:local-snapshot-20261012; GNU FDL 1.2 or later}
\begin{document}
\maketitle
\noindent\textit{Editorial note.} All delivered mathematical text is the author's own, copied verbatim from the pinned Stacks Project chapter file intersection.tex; no statement or proof step was rewritten, reordered or removed.\par\medskip
\noindent\textit{Editorial note.} Span recomputation by label pairing (LANE\_BRIEF 10.2): in this chapter the author's proof environment follows the statement environment as a standalone block; the spec covers the complete environment pair (statement span, proof span) as recorded in the pool candidate.\par\medskip
\noindent\textit{Editorial note.} Occupation scan: combined claim\_key + statement\_anchor + generic-line scan over all registration files (active and withdrawn) found 0 hits for this chapter (is-occupied.json, 2026-10-14); consistent with the duty manager's six-check dispatch.\par\medskip
\ifdefined\hypertarget\hypertarget{sourcepointer-1}{}\fi
\noindent\textit{Source reference.} Original reference 01XZ: sources/stacks-chow/tags.txt, source lines 2538--2538. Cross-chapter standard prerequisite cited by the author inside the counted span; it is used as a pointer and its own proof is not part of this item.\par\medskip
\ifdefined\hypertarget\hypertarget{sourcepointer-2}{}\fi
\noindent\textit{Source reference.} Original reference 08DX: sources/stacks-chow/tags.txt, source lines 10653--10653. Cross-chapter standard prerequisite cited by the author inside the counted span; it is used as a pointer and its own proof is not part of this item.\par\medskip
\ifdefined\hypertarget\hypertarget{sourcepointer-3}{}\fi
\noindent\textit{Source reference.} Original reference 0AZ4: sources/stacks-chow/tags.txt, source lines 13812--13812. Cross-chapter standard prerequisite cited by the author inside the counted span; it is used as a pointer and its own proof is not part of this item.\par\medskip
\ifdefined\hypertarget\hypertarget{sourcepointer-4}{}\fi
\noindent\textit{Source reference.} Original reference 08E8: sources/stacks-chow/tags.txt, source lines 10664--10664. Cross-chapter standard prerequisite cited by the author inside the counted span; it is used as a pointer and its own proof is not part of this item.\par\medskip
\ifdefined\hypertarget\hypertarget{sourcepointer-5}{}\fi
\noindent\textit{Source reference.} Original reference 09J3: sources/stacks-chow/tags.txt, source lines 12061--12061. Cross-chapter standard prerequisite cited by the author inside the counted span; it is used as a pointer and its own proof is not part of this item.\par\medskip
\begin{lemma}
\label{lemma-tensor-coherent}
Let $X$ be a locally Noetherian scheme.
\begin{enumerate}
\item If $\mathcal{F}$ and $\mathcal{G}$ are coherent $\mathcal{O}_X$-modules,
then $\text{Tor}_p^{\mathcal{O}_X}(\mathcal{F}, \mathcal{G})$ is too.
\item If $L$ and $K$ are in $D^-_{\textit{Coh}}(\mathcal{O}_X)$, then
so is $L \otimes_{\mathcal{O}_X}^\mathbf{L} K$.
\end{enumerate}
\end{lemma}

\begin{proof}
Let us explain how to prove (1) in a more elementary way and part (2)
using previously developed general theory.

\medskip\noindent
Proof of (1). Since formation of $\text{Tor}$ commutes with localization
we may assume $X$ is affine. Hence $X = \Spec(A)$ for some Noetherian
ring $A$ and $\mathcal{F}$, $\mathcal{G}$ correspond to finite $A$-modules
$M$ and $N$ (Cohomology of Schemes, Lemma
\ref{coherent-lemma-coherent-Noetherian}).
By Derived Categories of Schemes, Lemma
\ref{perfect-lemma-quasi-coherence-tensor-product} we may
compute the $\text{Tor}$'s by first computing the $\text{Tor}$'s
of $M$ and $N$ over $A$, and then taking the associated $\mathcal{O}_X$-module.
Since the modules $\text{Tor}_p^A(M, N)$ are finite by
Algebra, Lemma \ref{algebra-lemma-tor-noetherian}
we conclude.

\medskip\noindent
By Derived Categories of Schemes, Lemma
\ref{perfect-lemma-identify-pseudo-coherent-noetherian}
the assumption is equivalent to asking $L$ and $K$ to be
(locally) pseudo-coherent. Then $L \otimes_{\mathcal{O}_X}^\mathbf{L} K$
is pseudo-coherent by
Cohomology, Lemma \ref{cohomology-lemma-tensor-pseudo-coherent}.
\end{proof}

\end{document}
